\chapter{Démarche détaillée et implémentation actuelle}

Ce chapitre décrit la version effectivement présente dans le dépôt. Il complète les vues
fonctionnelles précédentes par une lecture reproductible des traitements, des fichiers
intermédiaires et des décisions de conception. La règle directrice est de séparer trois
responsabilités : construire la connaissance, sélectionner les éléments pertinents, puis
rédiger une réponse. Un LLM ne doit pas remplir simultanément ces trois rôles.

\section{Organisation actuelle du dépôt}

Le dépôt a été réorganisé afin de distinguer le code actif, les données, les essais et les
anciens prototypes. Le point d'entrée reste \code{api.py}; le code applicatif réutilisable
est regroupé dans le package \code{arianes/}.

\begin{lstlisting}[caption={Arborescence fonctionnelle actuelle}]
api.py                         API FastAPI et point d'entree
index.html                     interface web
arianes/
  agent.py                     orchestration et generation
  retrieval.py                 routage et retrieval GraphRAG
  paths.py                     chemins centralises
  evaluation.py                evaluation fonctionnelle
  graph/
    uml_builder.py             import UML et documents
    metier_builder.py          import de la couche metier
  pipelines/
    xmi_to_json.py             conversion XMI vers JSON
    extraction_metier.py       extraction d'une fiche
    extraction_globale.py      traitement du corpus
    validation_metier.py       validation V3
    recontrole_metier.py       second controle hors ligne
    selection_prudente.py      selection globale graduee
    generer_pilotes.py         pilote par document
    normalisation_metier.py    fusion de doublons
project_data/
  xmi/                         sources du metamodele
  uml/                         JSON structures
  procedures/                  fiches texte
  metier/                      propositions et pilotes
tests/                         tests automatises
ml/                            prototype de reranking
documentation/                 documentation technique
legacy/                        anciens prototypes non actifs
\end{lstlisting}

Le module \code{arianes/paths.py} centralise les chemins. Cette décision évite que chaque
script reconstruise des chemins relatifs différents et rend les commandes exécutables
depuis la racine. Le dossier \code{legacy/} est conservé pour la traçabilité, mais il ne
participe pas au lancement normal.

\section{Deux temporalités : hors ligne et en ligne}

\subsection{Préparation hors ligne}

La préparation hors ligne est exécutée lorsque les sources changent. Elle comprend la
conversion XMI, la construction du graphe UML, le découpage des fiches, le calcul des
embeddings, puis l'extraction et la validation du graphe métier. Ces opérations sont plus
coûteuses, mais leurs résultats sont persistés dans Neo4j et dans des fichiers JSON auditables.

\begin{center}
\resizebox{0.96\linewidth}{!}{%
\begin{tikzpicture}[node distance=7mm and 8mm]
\node[source] (xmi) {XMI Arianes};
\node[traitement,right=of xmi] (json) {Conversion\\\code{xmi\_to\_json.py}};
\node[stockage,right=of json] (uml) {JSON UML};
\node[traitement,below=of uml] (builder) {Construction\\UML et documents};
\node[stockage,left=of builder] (neo) {Neo4j};
\node[source,below=of xmi] (docs) {29 fiches texte};
\node[traitement,right=of docs] (extract) {Extraction et\\validation V3};
\node[stockage,right=of extract] (pilot) {Pilotes métier};
\draw[fleche] (xmi)--(json);
\draw[fleche] (json)--(uml);
\draw[fleche] (uml)--(builder);
\draw[fleche] (docs)--(builder);
\draw[fleche] (builder)--(neo);
\draw[fleche] (docs)--(extract);
\draw[fleche] (uml)--(extract);
\draw[fleche] (extract)--(pilot);
\draw[fleche] (pilot)--(neo);
\end{tikzpicture}
}
\end{center}

\subsection{Traitement en ligne}

Le traitement en ligne commence à la réception d'une question. Il ne reconstruit pas le
graphe. Il consulte les index et les relations déjà préparés, produit un contexte compact,
puis appelle le fournisseur choisi par l'utilisateur. La séparation garantit qu'un échec
du fournisseur de génération ne modifie jamais la connaissance stockée.

\section{Conversion XMI et construction UML}

\code{arianes/pipelines/xmi\_to\_json.py} analyse les exports XMI avec leur espace de noms.
Il extrait les classes, descriptions, attributs, multiplicités, énumérations, généralisations
et associations. Une étape de déduplication stabilise les listes lorsqu'un export contient
plusieurs représentations du même élément. Le JSON obtenu joue le rôle de contrat entre le
format externe XMI et Neo4j : le builder n'a pas besoin de connaître les détails XML.

\code{arianes/graph/uml\_builder.py} importe ensuite les objets par \code{MERGE}. Les
contraintes d'unicité portent notamment sur \code{Classe.nom}, \code{Attribut.key},
\code{Document.id} et \code{Chunk.id}. Une classe possède ses attributs par
\code{A\_ATTRIBUT}; un attribut pointe vers son type par \code{EST\_DE\_TYPE}. Les
relations entre classes sont limitées à \code{HERITE\_DE}, \code{ASSOCIE\_A},
\code{COMPOSE\_DE} et \code{AGREGE}.

Chaque fiche textuelle devient un \code{Document}. Elle est découpée par un séparateur de
texte avec chevauchement afin d'éviter de couper brutalement une information à la frontière
de deux passages. Chaque \code{Chunk} conserve son texte, sa position et son embedding.
Deux index complémentaires sont créés : un index plein texte pour les termes précis et un
index vectoriel pour les reformulations sémantiques.

\section{Routage explicite des questions}

\code{arianes/retrieval.py} commence par normaliser la question : suppression des accents,
passage en minuscules et découpage des noms CamelCase. Il résout ensuite les classes par
correspondance exacte, recouvrement lexical et similarité de chaîne. Les entités ambiguës
ne sont pas remplacées arbitrairement. Par exemple, \emph{Habillage} ne désigne aucune
classe exacte; le système propose alors \code{HabillageParcours}, \code{HabillageValeur}
ou d'autres classes réelles et demande une précision.

Le routage distingue actuellement six modes :

\begin{description}
  \item[STRUCTURE] consultation des classes, attributs et relations UML;
  \item[DOCUMENTAIRE] recherche de définitions, règles ou étapes dans les chunks;
  \item[METIER] interrogation des concepts et relations issus des fiches;
  \item[CHEMIN\_UML] plus court chemin multi-sauts entre deux classes UML;
  \item[CHEMIN\_METIER] plus court chemin entre deux concepts métier actifs;
  \item[HYBRIDE] combinaison ciblée de plusieurs couches lorsque la formulation le justifie.
\end{description}

Le qualificatif explicite de l'utilisateur est prioritaire. Ainsi, « chemin UML » ne doit
jamais être redirigé vers les relations métier. Cette distinction a été introduite après
l'observation d'un cas où \emph{Habillage} et \emph{Ligne} étaient recherchés dans la
mauvaise couche.

\section{Recherche des chemins multi-sauts}

Pour un chemin UML, deux classes exactes sont résolues avant la requête. Neo4j recherche
le plus court chemin, borné à huit relations, parmi les quatre types structurels autorisés.
Le parcours est non orienté pour découvrir une connexion, mais la réponse conserve le sens
réel de chaque relation. Lorsqu'une arête est parcourue à l'envers, ce fait est signalé.
Cette précaution évite de présenter une association inversée comme une nouvelle règle.

Pour un chemin métier, la profondeur est bornée à cinq relations. Seuls les concepts et
relations portant \code{valide=true} sont parcourus. Le résultat contient les preuves,
conditions, exceptions et statuts. Si aucune connexion n'existe, le système l'indique sans
demander au LLM de compléter le chemin.

\section{Recherche documentaire hybride}

La recherche vectorielle calcule l'embedding de la question puis compare ce vecteur aux
embeddings des chunks avec la similarité cosinus. En parallèle, l'index plein texte retrouve
les expressions exactes. Les deux classements sont fusionnés par Reciprocal Rank Fusion :
\[
  \mathrm{RRF}(d)=\sum_{r\in R}\frac{1}{60+\operatorname{rang}_{r}(d)}.
\]
Le score ne remplace pas les règles d'ancrage. Lorsqu'aucune classe n'est résolue, une
correspondance lexicale minimale est exigée pour éviter qu'une simple proximité vectorielle
ne ramène un passage hors sujet. Les chunks explicitement associés à une preuve métier sont
placés en priorité et dédupliqués.

Le prototype de reranking de \code{ml/} calcule des caractéristiques telles que score
sémantique, score lexical, score fusionné et recouvrement des mots. Il reste désactivé, car
le faible jeu d'entraînement ne démontre pas un gain par rapport à la fusion déterministe.

\section{Pipeline métier : de la fiche à Neo4j}

\subsection{Extraction contrôlée}

\code{extraction\_globale.py} découvre les fichiers \code{project\_data/procedures/*.txt}
et appelle \code{extraction\_metier.py} document par document. Des checkpoints permettent
de reprendre après interruption. Pour chaque morceau de texte, le programme sélectionne au
maximum trente classes UML candidates. La sélection est calculée avant l'appel au LLM;
elle limite le contexte et interdit au modèle de proposer librement une classe inexistante.

Le prompt demande un JSON comportant des concepts, leurs alias, leur catégorie, leurs
classes Arianes candidates et des triplets source--prédicat--cible. Chaque relation doit
être accompagnée d'une citation exacte. La sortie reste une \emph{proposition} : elle est
écrite dans \code{project\_data/metier/propositions/} et n'est pas importée automatiquement.

\subsection{Validation V3 graduée}

\code{validation\_metier.py} applique cinq critères pondérés : preuve exacte (30 points),
concepts identifiés (20), prédicat confirmé (20), direction confirmée (15) et condition
complète (15). Le score maximal est 100. Une erreur certaine, telle qu'une auto-relation,
une extrémité absente ou une preuve invérifiable, conduit à \code{REJETE}. En l'absence de
motif bloquant, un score d'au moins 85 avec passage disponible conduit à \code{VALIDE};
les autres relations exploitables sont classées \code{PROBABLE}.

\attention{Dans ce rapport, \code{VALIDE} signifie confirmé par les contrôles
documentaires. Il ne signifie pas certifié par un expert ferroviaire. Une relation
\code{PROBABLE} est utilisable seulement avec une formulation prudente et ses conditions.}

\code{recontrole\_metier.py} relit les propositions à partir des passages sources et écrit
une version V3 par document ainsi qu'une consolidation globale. \code{generer\_pilotes.py}
produit un pilote documentaire pour chaque fiche. \code{selection\_prudente.py} réunit les
relations \code{VALIDE} et \code{PROBABLE}, écarte les rejets et supprime les concepts non
utilisés. Le fichier \code{pilote\_global\_documentaire.json} est l'entrée de l'import global.

\subsection{Rôle des dossiers intermédiaires}

\begin{itemize}
  \item \code{propositions/} contient la sortie brute contrôlable du LLM;
  \item \code{corrections\_v2/} conserve l'ancien contrôle pour l'historique;
  \item \code{corrections\_v3/} contient le recontrôle gradué actuel;
  \item \code{pilotes\_documentaires/} contient une sélection exploitable par fiche;
  \item \code{pilote\_global\_documentaire.json} consolide le corpus pour l'import;
  \item \code{sauvegardes/} conserve l'état avant certaines fusions de doublons.
\end{itemize}

\subsection{Import et traçabilité}

\code{metier\_builder.py} refuse les relations rejetées, fusionne les concepts par identifiant
canonique et crée \code{RELATION\_METIER}. Les propriétés comprennent le statut, le score,
la condition, l'exception, les avertissements et les identifiants de chunks. Les ponts
\code{REPRESENTE\_PAR}, \code{DOCUMENTE\_PAR} et \code{ETAYE\_PAR} relient respectivement
un concept à une classe UML, un document et ses preuves textuelles. La couche métier reste
donc distincte du métamodèle tout en étant navigable dans la même base.

\section{Construction de la réponse}

\code{arianes/agent.py} ne décide pas quelles données lire : il reçoit le résultat du
retriever. Pour un chemin, une première représentation déterministe est produite afin de
figer les classes et relations. Pour une réponse générative, le contexte JSON contient
uniquement les éléments retenus. Le prompt interdit l'ajout d'une classe, d'une relation ou
d'une procédure absente et impose une formulation prudente pour \code{PROBABLE}.

L'utilisateur choisit explicitement Mistral, Gemini ou Groq dans l'interface. Il n'existe pas
de basculement silencieux entre fournisseurs : une erreur 429 ou 413 est associée au
fournisseur sélectionné. Ce choix facilite le diagnostic, car un résultat ne dépend pas d'un
fallback invisible. Les embeddings documentaires restent actuellement calculés avec le
service configuré lors de la construction et de la recherche.

\section{API, authentification et mémoire}

\code{api.py} initialise l'agent au démarrage, expose l'interface et stocke les comptes,
sessions, conversations et messages dans SQLite. Les mots de passe sont dérivés avec
PBKDF2-HMAC-SHA256, un sel et 600000 itérations. Le cookie de session est HttpOnly et
SameSite strict. Chaque accès à une conversation vérifie son propriétaire.

La mémoire conversationnelle n'est pas stockée dans Neo4j. SQLite conserve l'historique
par utilisateur; le retriever ne réutilise le message précédent que pour les relances
référentielles telles que « et ses attributs? ». Une nouvelle question autonome ne doit pas
être polluée par l'ancien sujet. Le graphe affiché dans l'interface provient du résultat de
retrieval de la question et constitue un complément visuel à la réponse textuelle.

\section{Reproductibilité et commandes}

\begin{lstlisting}[caption={Commandes principales depuis la racine}]
python -m arianes.pipelines.xmi_to_json
python -m arianes.graph.uml_builder
python -m arianes.pipelines.extraction_globale
python -m arianes.pipelines.recontrole_metier \
  --output-dir project_data/metier/corrections_v3
python -m arianes.pipelines.generer_pilotes
python -m arianes.pipelines.selection_prudente
python -m arianes.graph.metier_builder \
  project_data/metier/pilote_global_documentaire.json
python -m unittest discover -s tests -v
python -m uvicorn api:app --reload
\end{lstlisting}

La séparation des commandes permet de reconstruire seulement l'étape nécessaire et de
contrôler chaque artefact avant l'import. Elle facilite également l'explication du projet :
les scripts de préparation construisent une base traçable; le retriever sélectionne; le LLM
formule; l'API conserve la mémoire et applique les droits d'accès.
